\subsection{保留 KL 子问题与保留原教师目标的 MD 有别}

标准 entropy MD 对一个固定目标 $F$ 执行
\begin{equation}
 p^+=\argmin_p\{\langle\nabla F(q),p-q\rangle+\eta^{-1}\KL(p\Vert q)\}.
\end{equation}
取 $F_0(p)=\KL(p\Vert T)$，有 $\nabla F_0(q)=\log(q/T)+1$，与 $-r$ 只差常数，因而 $\kappa=0$ 的几何目标就是对原 reverse KL 的标准 entropy MD。

若 $\kappa>0$，有效梯度为 $-A_{\kappa,\eta}$，它本身依赖 $\eta$。其导数对 $r$ 为
\begin{equation}
 \frac{d}{dr}A_{\kappa,\eta}(e^r)=\frac{e^{\kappa r}}{1-\eta+\eta e^{\kappa r}}.
\end{equation}
当 $0<\eta<1$ 时它不是恒等于一的常数，所以一般不能用一个标量重新设定步长恢复原来的 $-r$ 方向。式~\eqref{eq:power-bpo-mirror-subproblem} 是精确的 entropy mirror 子问题，但不能称为对原 $F_0$ 的同一个 MD 更新。$\eta=1$ 或 $q=T$ 是特殊退化情况。信号依赖 $\eta$ 不妨碍为每个固定 $(\kappa,\eta)$ 构造另一个不依赖 anchor 的目标；本节已经给出新目标 $\mathcal J$ 的生成函数及收敛证明。

\subsection{连续极限是 teacher 散度的自然梯度流}

对 $\eta\to0$，记 $\rho_a=T_a/q_a$、$a_\kappa=\E_q\rho^\kappa$。逐项展开 $h_a=q_a[1+\eta(\rho_a^\kappa-1)/\kappa+O(\eta^2)]$ 并归一化，得到
\begin{equation}
 m_a=q_a+\eta v_{\kappa,a}+O(\eta^2),\qquad
 v_{\kappa,a}=\frac{q_a}{\kappa}(\rho_a^\kappa-a_\kappa).
 \label{eq:proximal-natural-velocity}
\end{equation}
展开针对固定完整支持的 $q,T$；其余项不能在任意极端概率比上直接视为统一 $O(\eta^2)$。$\kappa=0$ 的极限为 $v_{0,a}=q_a[\log\rho_a-\E_q\log\rho]$，$\kappa=1$ 为 $v_1=T-q$。

概率单纯形的 Fisher 度量为 $\langle v,w\rangle_q=\sum_av_aw_a/q_a$。若 $g_a$ 是 $L$ 的概率梯度代表，则负自然梯度为 $-q_a(g_a-\E_qg)$：其分量和为零，且对任何切向量 $w$，Fisher 内积恰好等于 $-\sum_ag_aw_a$。取 $L(q)=D_\kappa(T\Vert q)$，可取 $g_a=-\rho_a^\kappa/\kappa$；规范化生成函数给出的代表为 $(1-\rho_a^\kappa)/\kappa$，相差的常数被中心化消掉。因此 $v_\kappa=-\operatorname{grad}_{\!F}D_\kappa(T\Vert q)$。连续时间 $t$ 下
\begin{equation}
 \frac{dq_a}{dt}=v_{\kappa,a},\qquad
 \frac{d}{dt}D_\kappa(T\Vert q)
 =-\|v_\kappa\|_F^2
 =-\frac1{\kappa^2}\Var_q(\rho^\kappa).
 \label{eq:proximal-flow-lyapunov}
\end{equation}
最后一步由 $\sum_av_a=0$ 及 $g^\top v=-(\E_q\rho^{2\kappa}-a_\kappa^2)/\kappa^2$ 得到。$\kappa=0$ 时右侧为 $-\Var_q(\log\rho)$。所有形状参数都对应同一 Fisher 度量，但所下降的全局 teacher 散度不同。
\subsection{对固定幂散度执行有限步 MD 的另一种更新}

若要求在有限步长下严格采用这个固定 $F_\kappa$ 的 entropy MD，应另取
\begin{equation}
 \widetilde m_a\propto q_a
   \exp\left(\tau\frac{(T_a/q_a)^\kappa-1}{\kappa}\right).
 \label{eq:power-bpo-true-md-variant}
\end{equation}
它与幂均值目标只有一阶一致，有限 $\tau$ 时一般不同；并且指数中的概率比幂会重新引入极端正信号。另一种解释是第~\ref{sec:power-projection}~节 escort 坐标中的自适应二次目标梯度步，但这改变了坐标和目标，不应据此给实际神经网络 inner 更新贴上标准 MD 标签。

直接把 $D_\kappa(q\Vert p)$ 当成原负熵的 Bregman 正则也不正确：对 $0<\kappa<1$，它对 $p$ 的 Hessian 为 $\operatorname{diag}(q_a^\kappa p_a^{-\kappa-1})$，依赖参照 $q$。固定 mirror 势的 $B_\Phi(p,q)$ 对第一自变量的 Hessian 是 $\nabla^2\Phi(p)$，不依赖 $q$。一般情况下，两者不是同一 Bregman 散度。
